
\section{Methods}

\subsection{Pipeline Architecture}

GlycoSMILES2BAP operates through four sequential modules: Input Parsing, CCD Mapping, BAP Generation, and Output Formatting. The pipeline accepts IUPAC-condensed, WURCS, and GlycoCT input formats.

\subsection{CCD Mapper Module}

The CCD (Chemical Component Dictionary) provides standardized residue identifiers. Our mapper supports 28+ monosaccharide configurations. Key design decisions include: (1) case-insensitive matching, (2) anomeric position tracking (C2 for sialic acids, C1 for aldoses), and (3) ring oxygen positions (O4 for pentoses, O5 for hexoses, O6 for sialic acids).

\begin{table}[h]
\centering
\caption{Key CCD Mappings}
\begin{tabular}{lllll}
\toprule
Monosaccharide & Anomer & Config & CCD Code & Anomeric C \\
\midrule
GlcNAc & $\beta$ & D & NAG & C1 \\
GlcNAc & $\alpha$ & D & A2G & C1 \\
Man & $\alpha$ & D & MAN & C1 \\
Man & $\beta$ & D & BMA & C1 \\
Gal & $\beta$ & D & GAL & C1 \\
Gal & $\alpha$ & D & GLA & C1 \\
Fuc & $\alpha$ & L & FUC & C1 \\
Neu5Ac & $\alpha$ & D & SIA & C2 \\
Neu5Gc & $\alpha$ & D & NGC & C2 \\
Xyl & $\beta$ & D & XYS & C1 \\
\bottomrule
\end{tabular}
\end{table}

\subsection{BAP Generator Module}

The bondedAtomPairs (BAP) generator converts glycan topology into AF3-compatible bond specifications. Each glycosidic linkage is represented as an explicit atom pair specifying donor residue, donor atom, acceptor residue, and acceptor atom.

\subsection{Evaluation Metrics}

We define three accuracy metrics: Epimer Accuracy measures correct preservation of monosaccharide stereochemistry; Anomeric Accuracy measures correct $\alpha/\beta$ configuration; Linkage Accuracy measures correct donor and acceptor atom positions.

